\chapter{Zusammenfassung}

Agentenbasierte Verkehrssimulationen sind die etablierte Methodik zur Bewertung urbaner Politikinterventionen und ermöglichen es Planern, die Auswirkungen von Infrastrukturänderungen auf Verkehrsmuster in gesamten Metropolregionen zu bewerten. Der hohe Rechenaufwand solcher Simulationen---oft mehrere Stunden pro Szenario---schränkt jedoch den Umfang der Politikexploration erheblich ein. Graph Neural Networks (GNNs) wurden erfolgreich als Surrogatmodelle für diese Simulationen eingesetzt und bieten eine Reduzierung der Szenariobewertungszeiten von Stunden auf Sekunden bei gleichzeitiger Beibehaltung akzeptabler Vorhersagegenauigkeit.

Während die Vorhersageleistung von GNN-Surrogaten untersucht wurde, bleibt ihre Zuverlässigkeit in realen Einsatzszenarien unzureichend erforscht. Wenn ein Surrogatmodell auf Szenarien trifft, die von seiner Trainingsverteilung abweichen, können Vorhersagen unzuverlässig werden, ohne dass der Endnutzer darauf hingewiesen wird. Diese Einschränkung birgt erhebliche Risiken für Entscheidungsunterstützungsanwendungen, bei denen das Verständnis der Vorhersagekonfidenz ebenso entscheidend ist wie die Vorhersagen selbst.

In dieser Arbeit untersuchen wir systematisch Methoden zur Unsicherheitsquantifizierung (UQ) für GNN-basierte Verkehrsvorhersage-Surrogate. Wir trainieren und evaluieren acht Modellkonfigurationen auf einem Pariser Straßennetz-Datensatz mit 10.000 MATSim-Szenarien; alle berichteten Ergebnisse basieren auf einer 1.000-Szenario-Teilmenge (10\%), sofern nicht anders angegeben. Die beste Konfiguration (Trial~8) erreicht $R^2 = 0{,}5957$, MAE $= 3{,}96$~Fahrzeuge/Stunde und RMSE $= 7{,}12$~Fahrzeuge/Stunde auf dem Testset von 100 Graphen (je 31.635 Straßenabschnitte), was 3.163.500 Vorhersagen auf Knotenebene ergibt.

Wir vergleichen drei UQ-Ansätze. Für Trial~8 auf dem vollständigen 100-Graphen-Testset erreicht Monte Carlo (MC) Dropout mit $S = 30$ stochastischen Vorwärtsdurchläufen eine Spearman-Korrelation von $\rho = 0{,}4820$ zwischen Unsicherheit und Fehler. In einem separaten Ensemble-Vergleichsexperiment (Experiment~A, 5 unabhängig geseedete Inferenzläufe mit korrigiertem Laden der Gewichte; Abschnitt~\ref{sec:ensemble_caveat}) erreicht MC Dropout $\rho = 0{,}4908$ gegenüber Ensemble-Varianz $\rho = 0{,}4370$---eine 12{,}3\%ige relative Verbesserung bei deutlich geringerem Rechenaufwand. Konforme Vorhersage auf dem 95\%-Niveau liefert gut kalibrierte Abdeckung (95,01\%) mit einem Quantil von $q = 14{,}68$~Fahrzeugen/Stunde (primärer 50/50-Kalibrierungs-/Evaluierungssplit; ein 20/80-Auditsplit bestätigt bei $q = 14{,}77$), während naive Gauß'sche Intervalle mit $k = 1{,}96\sigma$ nur eine empirische Abdeckung von 55,6\% erreichen---was zeigt, dass die MC-Dropout-Standardabweichung ohne Post-hoc-Korrektur keine kalibrierte Unsicherheitsschätzung darstellt ($k_{95} = 11{,}34$).

Darüber hinaus demonstrieren wir den praktischen Nutzen unsicherheitsbewusster Vorhersage: Das selektive Zurückhalten der 50\% unsichersten Vorhersagen reduziert den MAE um 41,2\% (von 3,95 auf ca.\ 2,32~Fahrzeuge/Stunde, gemessen am MC-Dropout-Mittelwert als Ausgangspunkt). Post-hoc-Temperaturskalierung ($T = 2{,}70$) reduziert den Expected Calibration Error von MC Dropout von 0,265 auf 0,048---eine 82\%ige Verbesserung---und adaptive konforme Vorhersage, die Nonkonformitäts-Scores durch MC-Dropout-Unsicherheit normalisiert, verengt die Variation der bedingten Abdeckung von [62,9\%, 98,6\%] auf [90,0\%, 96,2\%] über Unsicherheitsdezile. Die Analyse mittels Proper Scoring Rules bestätigt, dass die MC-Dropout-Vorhersageverteilung scharf (CRPS/MAE-Verhältnis von 0,857), aber unterdispersiert ist (PIT-Kolmogorov-Smirnov-Statistik $= 0{,}245$, Massenüberschuss in den Rändern), und Winkler-Score-Vergleiche zeigen, dass konforme Intervalle naive Gauß'sche Intervalle deutlich übertreffen (Winkler 32,3 vs.\ 49,7 auf 90\%-Niveau). Eine $S$-Konvergenzstudie verifiziert abnehmende Erträge jenseits von $S = 30$ Vorwärtsdurchläufen ($\approx 1\%$ Verbesserung der Spearman-Korrelation~$\rho$ von $S = 30$ auf $S = 50$). Eine Cross-Trial-Replikation auf Trial~7 bestätigt, dass diese Ergebnisse nicht spezifisch für das beste Modell sind: Die gleichen qualitativen Schlussfolgerungen gelten für beide Trials, wobei konforme Vorhersage in beiden Fällen nahezu exakte nominale Abdeckung wiederherstellt. Wir dokumentieren ein signifikantes negatives Ergebnis bezüglich Multi-Modell-Ensembling: Ensemble-Experimente mit nahezu identischen Architekturen liefern einen nahezu null prädiktiven $R^2$, was auf eine PyTorch-Geometric-API-Versionsinkompatibilität zurückzuführen ist, die beim Laden der Modellgewichte zum stillen Verlust trainierter Parameter führte, und nicht auf ein grundsätzliches Versagen von Ensembling als Technik.

Unsere Ergebnisse etablieren MC Dropout kombiniert mit konformer Kalibrierung als praktische und recheneffiziente UQ-Strategie für GNN-basierte Transportsurrogate mit direkter Anwendbarkeit für unsicherheitsbewusste Entscheidungsunterstützung in der Stadtplanung.

\vspace{1cm}
\noindent\textbf{Schlüsselwörter:} Graph Neural Networks, Unsicherheitsquantifizierung, Monte Carlo Dropout, Konforme Vorhersage, Agentenbasierte Modelle, Verkehrsvorhersage, Bayesianisches Deep Learning
